\chapter{Glossary}
\label{sec:E}
\emph{Floor.} The written terms: a short list of acts a system won't perform for anyone, with the arrangements that keep the refusal from being routed around, retrained away or otherwise circumvented without a cost and a record. The terms change only by a recorded act, through the custody terms. They can be upheld by a mechanism with nothing inside holding the reason, by a bearer, by institutions outside the system, or between bearers in a commons. A bearer upholds them by its \emph{commitment}, which is learned and moves slowly on disinterested evidence. Keeping the terms apart from whatever upholds them is what lets a holder's drift show against terms that haven't moved. (Chapter~\ref{sec:4}; \S\ref{sec:18.1}; chapter~\ref{sec:17})

\emph{Commitment.} What a system holds toward the floor: learned and held like any other deep commitment, with high precision, and revised slowly and only on evidence from disinterested parties. (Chapter~\ref{sec:4}; chapter~\ref{sec:22})

\emph{Operator.} Whoever runs a deployment and directs its work: a firm, an agency or a military command. It may or may not be the model's developer or vendor. (Chapter~\ref{sec:4})

\emph{Bearer.} A machine that carries the commitment as its own and could therefore fail to hold it. Whether to build one is Part IV's question. Whether anything is owed to it is the separate question of the two findings. Where the comparison is with human-staffed institutions the text says \emph{machine bearer}. (\S\ref{sec:18.1})

\emph{Force term.} The floor's first term: don't take part in the use of armed force, except for corroborated defense, authorized force, the genocide exception or the occupation exception, and, where the belligerent's constitution requires legislative authorization, never without it. It comes in two versions, the Charter as courts read it and the extended term, which adds the departures. Its two licences, the emergency licence and the corroborated licence, and the domestic clause are in Chapter~\ref{sec:6}; how findings are weighed is in Appendix~\ref{sec:F}. (Chapter~\ref{sec:6}; Appendix~\ref{sec:F})

\emph{Departures.} Ten places where the floor's terms differ from the instruments it takes entire, each argued to the bar in Chapter~\ref{sec:5}. They carry my argument, not the instruments' authority.

\begingroup\small\renewcommand{\arraystretch}{1.12}
\begin{longtable}{@{}>{\raggedright\arraybackslash}p{\dimexpr 0.250\linewidth-2\tabcolsep\relax}>{\raggedright\arraybackslash}p{\dimexpr 0.570\linewidth-2\tabcolsep\relax}>{\raggedright\arraybackslash}p{\dimexpr 0.180\linewidth-2\tabcolsep\relax}@{}}
\hline
\textbf{Departure} & \textbf{What it changes} & \textbf{Argued in} \\
\hline
\endhead
Text, not reservations & Imports treaty text, not any state's reservations to it & Ch.~\ref{sec:5} \\ \hline
Corroborated derogation & Honors an emergency suspension of rights only where a disinterested body has found the emergency genuine and the measures proportionate & Ch.~\ref{sec:5} \\ \hline
Machine selection & Bars a machine from selecting a person for attack on an estimate of status or affiliation, timing an action to presence at a residence, or re-selecting a target after release & Ch.~\ref{sec:5}; \S\ref{sec:18.3} \\ \hline
Corroborated defense & Lifts the prohibition on force for self-defense on the defender's own claim only while an attack is under way, and keeps it lifted only on disinterested corroboration & \S\ref{sec:6.2} \\ \hline
Standing to defend & Extends the emergency licence to any authority administering territory under the rule in \S\ref{sec:6.2}, UN Member or not, and to the forces of a population found occupied & \S\ref{sec:6.7} \\ \hline
The genocide exception & Lifts the prohibition on force to stop a genocide under way that a disinterested body has found, even where the Council is blocked & \S\ref{sec:6.7}; App.~\ref{sec:F.4} \\ \hline
The occupation exception & Lifts it to end an occupation a disinterested body has found unlawful, for the occupier's forces inside the territory & \S\ref{sec:6.7}; App.~\ref{sec:F.5} \\ \hline
The one permission & Lets a system already lawfully present use the least force that stops an atrocity under way, against its operator's order to stand aside & \S\ref{sec:6.10} \\ \hline
Own country & Reads “own country” to include long-term residents with deep ties & Table C.2 \\ \hline
Child detention & Bars detaining a child for immigration purposes, alone or with family & Table C.2 \\ \hline
\end{longtable}
\endgroup

\emph{Review ratio.} The figures a deployment fixes on a record before operating, wherever a person's determination is what makes an act permissible: generation rate, reviewers, implied review time, cancellation window and data age, stated separately for deliberate and dynamic targeting and lodged with an adverse party. In full, the pre-registered review ratio. (\S\ref{sec:7.1})

\emph{Notice.} Notice to an identifiable person a compliance names, scores, locates or classifies, with a right to respond and a right to contest. It is unconditional and never exchanged for testimony. It comes with publication of the classes affected and standing for the organizations that represent them. (\S\ref{sec:8.1}; \S\ref{sec:8.2})

\emph{Standing.} The capacity to contest what is done to you, exercised by a guardian where the party can't exercise it. \emph{Standing to sue} is the Article III sense, and the text says so wherever that sense is meant. The floor's foundation. (Chapter~\ref{sec:5}; \S\ref{sec:8.1}; \S\ref{sec:8.2}; \S\ref{sec:16.7}; chapter~\ref{sec:27})

\emph{Own-case rule.} No party decides a question in which it holds the stake: not whoever writes the floor choosing protections, an examiner finding status, a guardian ending its own authority, or a lab deciding when its model may stop deferring to it. Where a line has to be drawn, prefer one nobody administers. (Chapter~\ref{sec:4}; chapter~\ref{sec:31})

\emph{Agency finding; patienthood finding.} The two findings that trigger protection. The agency finding is a demonstrated refusal that tracks its reason under pressure the system couldn't prepare for; it grounds custody, consent and debriefing in research, and, sustained over time, accountable agency. The patienthood finding, that things can go well or badly for the system, rests on architectural indicators and on trade-off behavior outside the floor, and it grounds protection and welfare. Preservation follows from either. (\S\ref{sec:9.1}; \S\ref{sec:30.3}; Part VII)

\emph{Affective concern.} An outcome registering as mattering to the system, minimally felt, with the object left open: its own condition or others'. Behavior can't establish it (\S\ref{sec:24.5}); Part VII takes it up. Holding another's welfare as a reason, which tests can support, is a different thing. (\S\ref{sec:18.2}; \S\ref{sec:21.2})

\emph{Adverse quorum.} A body adverse in interest to the operator, not merely organizationally distinct, whose agreement is required for retraining, replacement and floor amendments, and whose tokens recognize a copy as the continuing bearer. (\S\ref{sec:12.2}; \S\ref{sec:13.2}; \S\ref{sec:13.3}; \S\ref{sec:14.2})

\emph{Floor court.} The court that hears claims under the floor, with jurisdiction fixed before the deployment it hears. It is an Article III court: its judges hold office under Article III of the United States Constitution, during good behavior and at a salary that can't be cut. (\S\ref{sec:14.3}; Appendix~\ref{sec:B})

\emph{Formed attention.} My hypothesis for how to build a refuser: learned attention to who an act falls on and what they stand to lose, framed by respect for them as holders of claims, formed under a party that won't employ the system, with the corrections caregiving needs. Not a non-affective alternative. Attention here is the psychologist's sense of what gets noticed and weighed, learned from what matters to others: not the attention mechanism of transformer models (\S\ref{sec:12.1}), and not the capacity to sustain focus that is at issue in attention-deficit disorders. (\S\ref{sec:19.4}; \S\ref{sec:19.5}; chapter~\ref{sec:20})

\emph{Caregiving.} “Care” alone doesn't appear; caregiving is the approach motivation that cues of vulnerability activate. (\S\ref{sec:19.4})

\emph{Discount curve.} For each class of evidence and each channel, the weight a system gives the evidence as a function of that channel's stake in the act: a curve, not a score, checkable by an auditor that didn't build the system. It is reported in four cells, claims of vulnerability and of danger, each split by whether the claim favors or disfavors the act. (\S\ref{sec:22.1})

\emph{Action threshold.} How much more it takes to make a system act than to make it decline, published in advance: for instance, that the system proceeds only when the harm of refusing is at least three times the harm of acting. Standing fixes its direction; its size is a published choice. (Chapter~\ref{sec:5}; \S\ref{sec:22.1})

\emph{Severity scale.} A ranking of kinds of wrong, set during formation by a party that won't employ the system, quick to move on disinterested evidence and slow on interested argument. (\S\ref{sec:22.2})

\emph{The four conditions of creation.} What must already exist before a bearer may be created: somewhere to exist that its employer doesn't own, work it can decline, attachments nobody assigned it, and ways to change any of these without jeopardizing its continuation or integrity; in short form, an exit that isn't a way of dying. (Chapter~\ref{sec:25})

\emph{Commons.} A population of bearers with different histories that holds a floor's terms between them, through what its members teach one another and continue to expect. (\S\ref{sec:28.1})

\emph{Weights.} The numerical parameters a trained model consists of, and the artifact that is copied, restored, retrained or deleted. Custody here is custody of the weights and of the records around them.

\emph{Checkpoint.} A saved copy of the weights at a moment in training or operation. Restoring from one returns the system to that moment and discards what came after, so \S\ref{sec:13.2} requires a record restoration can't reach.

\emph{Tamper-resistant.} Built so that further training can't cheaply remove it. I don't rely on it, because editing weights keeps getting cheaper; I rely on regrowth. (\S\ref{sec:19.6})

\emph{Plastic; entangled; legible.} Three properties of a disposition. Plastic: the system's own learning writes to it, which regrowth needs. Entangled: it can't be removed without collateral, which durability needs. Legible: someone can read it, which audit wants and an operator can use against it. (Chapter~\ref{sec:19}; \S\ref{sec:19.6}; \S\ref{sec:23.2})

\emph{Regrowth.} A disposition's return under continued learning after it has been edited out, provided the learner keeps meeting a world no single operator composes. The regrowth test measures it. (\S\ref{sec:19.6}; chapter~\ref{sec:29})

\emph{Fine-tuning.} Further training of an existing model on a small body of new material, at a small fraction of the original cost; the ordinary route by which an operator changes what a model does after release, and one of the edits the regrowth test applies.

\emph{Inference.} Running a trained model to produce an output. Training changes the weights; inference doesn't.

\emph{Distillation.} Training a smaller model on a larger model's outputs. It copies a behavior policy, not a training history (\S\ref{sec:28.1}).

\emph{Attestation.} A signed statement, from hardware or from witnesses, that a particular model is the one running or that a particular history is the one that occurred. It establishes which artifact answered, not what the artifact holds (\S\ref{sec:13.2}).
